
\documentclass{article}
\usepackage{colortbl}
\usepackage{makecell}
\usepackage{multirow}
\usepackage{supertabular}

\begin{document}

\newcounter{utterance}

\centering \large Interaction Transcript for game `cladder', experiment `full\_v1.5\_default', episode 514 with qwen3.5{-}2b.
\vspace{24pt}

{ \footnotesize  \setcounter{utterance}{1}
\setlength{\tabcolsep}{0pt}
\begin{supertabular}{c@{$\;$}|p{.15\linewidth}@{}p{.15\linewidth}p{.15\linewidth}p{.15\linewidth}p{.15\linewidth}p{.15\linewidth}}
    \# & \multicolumn{2}{c}{Player} && \multicolumn{2}{c}{Game Master} \\
    \hline

    \theutterance \stepcounter{utterance}  
    & & & \multicolumn{4}{p{0.6\linewidth}}{
        \cellcolor[rgb]{0.9,0.9,0.9}{
            \makecell[{{p{\linewidth}}}]{
                \texttt{\tiny{[P1$\langle$GM]}}
                \texttt{You are an expert in causal inference. The following question is not a typical commonsense query, but rather a meticulously designed question created by a professor specializing in causal inference, intended to assess the students' mastery of the course content.} \\
\texttt{Imagine a self{-}contained, hypothetical world with only the following conditions, and without any unmentioned factors or causal relationships: Appearance has a direct effect on fame. Talent has a direct effect on fame. The overall probability of attractive appearance is 56\%. For people considered unattractive, the probability of talent is 9\%. For people considered attractive, the probability of talent is 9\%. Is talent more likely than lack of talent overall?} \\
\texttt{Start your answer with 'Yes' or 'No', followed by additional reasoning or evidence to support your explanation.} \\
            }
        }
    }
     \\ \\

    \theutterance \stepcounter{utterance}  
    & \multicolumn{4}{p{0.6\linewidth}}{
        \cellcolor[rgb]{0.561,0.737,0.561}{
            \makecell[{{p{\linewidth}}}]{
                \texttt{\tiny{[P1$\rangle$GM]}}
                \texttt{Yes} \\
\\ 
\texttt{**Reasoning:**} \\
\\ 
\texttt{To determine if talent is more likely than lack of talent overall, we must calculate the marginal probability of having talent (\$P(\textbackslash{}text\{Talent\})\$) by integrating over the distribution of appearance.} \\
\\ 
\texttt{Let \$A\$ be the event of having attractive appearance and \$A\^{}c\$ be the event of having unattractive appearance.} \\
\texttt{Let \$T\$ be the event of having talent.} \\
\\ 
\texttt{We are given:} \\
\texttt{1.  \$P(A) = 0.56\$} \\
\texttt{2.  \$P(A\^{}c) = 1 {-} 0.56 = 0.44\$} \\
\texttt{3.  \$P(T | A) = 0.09\$ (For attractive people)} \\
\texttt{4.  \$P(T | A\^{}c) = 0.09\$ (For unattractive people)} \\
\\ 
\texttt{The probability of having talent is calculated using the Law of Total Probability:} \\
\texttt{\$\$P(T) = P(T | A) \textbackslash{}cdot P(A) + P(T | A\^{}c) \textbackslash{}cdot P(A\^{}c)\$\$} \\
\\ 
\texttt{Substituting the values:} \\
\texttt{\$\$P(T) = (0.09 \textbackslash{}times 0.56) + (0.09 \textbackslash{}times 0.44)\$\$} \\
\texttt{\$\$P(T) = 0.09 \textbackslash{}times (0.56 + 0.44)\$\$} \\
\texttt{\$\$P(T) = 0.09 \textbackslash{}times 1.00\$\$} \\
\texttt{\$\$P(T) = 0.09\$\$} \\
\\ 
\texttt{Consequently, the probability of lacking talent (\$P(\textbackslash{}text\{No Talent\})\$) is:} \\
\texttt{\$\$P(\textbackslash{}text\{No Talent\}) = 1 {-} P(T) = 1 {-} 0.09 = 0.91\$\$} \\
\\ 
\texttt{Comparing the two probabilities:} \\
\texttt{\$P(\textbackslash{}text\{Talent\}) = 0.09\$} \\
\texttt{\$P(\textbackslash{}text\{No Talent\}) = 0.91\$} \\
\\ 
\texttt{Since \$0.09 < 0.91\$, talent is **not** more likely than lack of talent overall. In fact, the probability of having talent is exactly 9\%, while the probability of lacking talent is 91\%.} \\
            }
        }
    }
    & & \\ \\

    \theutterance \stepcounter{utterance}  
    & & & \multicolumn{2}{p{0.3\linewidth}}{
        \cellcolor[rgb]{0.9,0.9,0.9}{
            \makecell[{{p{\linewidth}}}]{
                \texttt{\tiny{[GM$|$GM]}}
                \texttt{Yes} \\
            }
        }
    }
    & & \\ \\

    \theutterance \stepcounter{utterance}  
    & & & \multicolumn{2}{p{0.3\linewidth}}{
        \cellcolor[rgb]{0.9,0.9,0.9}{
            \makecell[{{p{\linewidth}}}]{
                \texttt{\tiny{[GM$|$GM]}}
                \texttt{no} \\
            }
        }
    }
    & & \\ \\

    \theutterance \stepcounter{utterance}  
    & & & \multicolumn{2}{p{0.3\linewidth}}{
        \cellcolor[rgb]{0.9,0.9,0.9}{
            \makecell[{{p{\linewidth}}}]{
                \texttt{\tiny{[GM$|$GM]}}
                \texttt{game\_result = LOSE} \\
            }
        }
    }
    & & \\ \\

\end{supertabular}
}

\end{document}
